\documentclass[conference]{IEEEtran}
\IEEEoverridecommandlockouts

\usepackage{amsmath,amssymb,amsfonts}
\usepackage{algorithmic}
\usepackage{algorithm}
\usepackage{graphicx}
\usepackage{booktabs}
\usepackage{xcolor}
\usepackage{hyperref}
\usepackage{multirow}
\usepackage{bm}

\def\BibTeX{{\rm B\kern-.05em{\sc i\kern-.025em b}\kern-.08em
    T\kern-.1667em\lower.7ex\hbox{E}\kern-.125emX}}

\begin{document}

\title{RC-FedBiT: Rate-Controlled Federated Learning\\
with Bit-width Adaptation}

\author{
\IEEEauthorblockN{Maolin Wang}
\IEEEauthorblockA{
Hong Kong Institute of AI for Science (HKAI-Sci)\\
City University of Hong Kong\\
Hong Kong, China\\
morin.wang@cityu.edu.hk}
}

\maketitle

\begin{abstract}
Federated learning (FL) over wireless networks faces two compounding challenges: prohibitive communication costs incurred by transmitting full-precision gradients, and time-varying Rayleigh fading channels that degrade gradient fidelity unpredictably. We present \textbf{RC-FedBiT}, a rate-controlled federated learning framework that adapts gradient bit-width to instantaneous wireless channel conditions. RC-FedBiT integrates four mechanisms: (1)~a \emph{Rank-1 Gradient Compressor} that decomposes each layer update $\Delta\mathbf{W}$ into a 1-bit sign matrix $\mathbf{B}$ and two floating-point scaling vectors $(\mathbf{h}_1, \mathbf{h}_2)$, achieving up to $32\times$ compression with bounded Frobenius-norm error; (2)~\emph{Channel-Adaptive Rate Selection} (CA-RS) that dynamically selects payload precision (FP32 / INT8 / binary-only) based on per-client instantaneous SNR and cosine-annealed quality thresholds; (3)~\emph{Noise-Impaired Aggregation with Channel Variance Adaptation} (NIA-CVA) that reweights client contributions by reconstruction error and channel quality; and (4)~\emph{Adaptive Frequency Calibration} (AFC) that corrects systematic spectral drift introduced by binary quantization. Experiments on CIFAR-10 with ViT-Tiny ($N=100$ clients, $\rho=0.1$, $T=100$ rounds, $\overline{\gamma}=10$\,dB) demonstrate that RC-FedBiT achieves $26.29\%$ test accuracy using only $\sim\!25\%$ of FedAvg's communication budget, outperforming all low-communication baselines. Ablation variants reach $27.13\%$, surpassing full-precision FedAvg ($26.54\%$).
\end{abstract}

\begin{IEEEkeywords}
Federated learning, gradient compression, wireless communication, Rayleigh fading, bit-width adaptation
\end{IEEEkeywords}

%======================================================
\section{Introduction}
%======================================================

Federated learning~\cite{mcmahan2017communication} enables collaborative model training across $N$ distributed clients without centralizing raw data. In each communication round $t$, a subset of $K = \lceil\rho N\rceil$ clients download the global model $\mathbf{W}^{(t)}$, perform $E$ steps of local SGD, and upload gradient updates $\Delta\mathbf{W}_i$ to a parameter server. For modern neural networks, a single FP32 upload can exceed hundreds of megabytes per client per round, constituting a severe bottleneck in bandwidth-constrained wireless deployments such as mobile edge computing, vehicular FL, and IoT sensor networks.

The problem is compounded by wireless channel dynamics. In Rayleigh fading environments, the complex channel gain $h_i \sim \mathcal{CN}(0,1)$ fluctuates independently across clients and rounds. The instantaneous SNR $\gamma_i = |h_i|^2 P / \sigma_n^2$ follows an exponential distribution; at mean SNR $\overline{\gamma} = 10$\,dB, approximately $27\%$ of clients experience SNR below $5$\,dB in any given round, severely degrading high-precision gradient transmissions.

Existing compression methods lack channel awareness. Sign-based methods such as SignSGD~\cite{bernstein2018signsgd} reduce communication cost to $\sim\!3\%$ of FedAvg but discard all gradient magnitude information, incurring a $3.28$ percentage-point (pp) accuracy loss on our benchmark. Stochastic quantization methods such as QSGD~\cite{alistarh2017qsgd} achieve unbiased compression at $\sim\!25\%$ cost but ignore channel state; a client whose FP32 transmission is corrupted by a $0$\,dB channel provides worse updates than simply transmitting only the sign. Low-rank methods such as PowerSGD~\cite{vogels2019powersgd} use rank-$r$ approximation at $\sim\!30\%$ cost but are likewise channel-oblivious.

RC-FedBiT fills this gap by making four contributions:
\begin{enumerate}
    \item A \textbf{Rank-1 Gradient Compressor} with GPU-accelerated power iteration, achieving $\sim\!25\times$ compression with bounded approximation error.
    \item \textbf{Channel-Adaptive Rate Selection (CA-RS)} with cosine-annealed dynamic thresholds, selecting among FP32, INT8, and binary-only payloads per client per round.
    \item \textbf{NIA-CVA} aggregation weighting that suppresses high-noise client contributions based on reconstruction error and channel quality.
    \item \textbf{Adaptive Frequency Calibration (AFC)} that corrects binary quantization spectral bias accumulated over training rounds.
\end{enumerate}

%======================================================
\section{Problem Formulation}
%======================================================

\subsection{Federated Optimization}

We consider a federated system with $N$ clients. Client $i$ holds local dataset $\mathcal{D}_i$ with $n_i = |\mathcal{D}_i|$ samples and weight $p_i = n_i / \sum_j n_j$. The global objective is:
\begin{equation}
    \min_{\mathbf{W}}\; F(\mathbf{W}) = \sum_{i=1}^{N} p_i\, F_i(\mathbf{W}),
    \quad
    F_i(\mathbf{W}) = \mathbb{E}_{(\mathbf{x},y)\sim\mathcal{D}_i}\!\left[\ell(\mathbf{W};\mathbf{x},y)\right].
\end{equation}
In each round, each participating client $i$ performs $E$ local SGD steps from the global weights:
\begin{equation}
    \mathbf{w}_i^{(t,e+1)} = \mathbf{w}_i^{(t,e)} - \eta_\ell\,\nabla F_i\!\left(\mathbf{w}_i^{(t,e)}\right), \quad e = 0,\ldots,E-1,
\end{equation}
and computes the pseudo-gradient $\Delta\mathbf{W}_i^{(t)} = \mathbf{w}_i^{(t,E)} - \mathbf{W}^{(t)} \in \mathbb{R}^{m\times n}$.

\subsection{Wireless Channel Model}

Client $i$ transmits its compressed gradient over an independent Rayleigh fading channel:
\begin{equation}
    y_i = h_i x_i + n_i, \quad h_i \sim \mathcal{CN}(0,1),\; n_i \sim \mathcal{CN}(0,\sigma_n^2).
\end{equation}
The instantaneous SNR is $\gamma_i = |h_i|^2 P / \sigma_n^2$. Since $|h_i|^2 \sim \mathrm{Exp}(1)$:
\begin{equation}
    \gamma_i \sim \mathrm{Exp}(\overline{\gamma}), \qquad \overline{\gamma} = 10^{\,\overline{\gamma}_{\mathrm{dB}}/10},
\end{equation}
with CDF $P(\gamma_i \leq x) = 1 - e^{-x/\overline{\gamma}}$. At $\overline{\gamma} = 10$\,dB: $P(\gamma < 5\,\text{dB}) = 27.3\%$.

\subsection{Communication-Constrained Problem}

RC-FedBiT solves:
\begin{equation}
    \min_{\mathbf{W}} F(\mathbf{W}) \quad \text{s.t.} \quad \mathbb{E}\!\left[B_i^{(t)}\right] \leq B_{\max} \;\forall\, i, t,
\end{equation}
where $B_i^{(t)}$ is the bit count transmitted by client $i$ in round $t$.

%======================================================
\section{Method}
%======================================================

\subsection{Rank-1 Gradient Compression}

For each layer $\ell$ with weight update $\Delta\mathbf{W} \in \mathbb{R}^{m\times n}$, the Eckart–Young theorem gives the best rank-1 approximation:
\begin{equation}
    \Delta\mathbf{W} \approx \sigma_1\, \mathbf{u}_1 \mathbf{v}_1^\top,
\end{equation}
where $\sigma_1, \mathbf{u}_1, \mathbf{v}_1$ are the leading singular triplet. RC-FedBiT factors this as:
\begin{align}
    \mathbf{B} &= \mathrm{sign}(\Delta\mathbf{W}) \in \{-1,+1\}^{m\times n}, \\
    \mathbf{h}_1 &= \sqrt{\sigma_1}\,\mathbf{u}_1 \in \mathbb{R}^m, \quad \mathbf{h}_2 = \sqrt{\sigma_1}\,\mathbf{v}_1 \in \mathbb{R}^n, \\
    \widehat{\Delta\mathbf{W}} &= \mathbf{B} \odot \left(\mathbf{h}_1 \mathbf{h}_2^\top\right), \label{eq:recon}
\end{align}
where $\odot$ denotes element-wise multiplication. The Frobenius approximation error is:
\begin{equation}
    \left\|\Delta\mathbf{W} - \widehat{\Delta\mathbf{W}}\right\|_F^2 = \sum_{k=2}^{\min(m,n)} \sigma_k^2 \;\leq\; \|\Delta\mathbf{W}\|_F^2 - \sigma_1^2.
\end{equation}
The communication cost per layer is:
\begin{align}
    C_{\mathrm{FP32}} &= 32mn \text{ bits}, \\
    C_{\mathrm{RC}} &= mn + 32(m+n) \text{ bits}, \\
    R_{\mathrm{comp}} &= \frac{32mn}{mn + 32(m+n)} \;\xrightarrow{m=n=d}\; \frac{32d}{d+64}.
\end{align}
For $d = 192$ (ViT-Tiny attention): $R_{\mathrm{comp}} = 24\times$; for $d = 768$ (FFN): $R_{\mathrm{comp}} = 29\times$.

\textbf{GPU Power Iteration.} The leading singular triplet is computed via 3-step power iteration on GPU with $O(mn)$ complexity, compared to $O(mn\min(m,n))$ for full SVD.

\subsection{Channel-Adaptive Rate Selection (CA-RS)}

\textbf{SNR sampling.} At each round, the server samples client channel SNRs via inverse CDF:
\begin{equation}
    \gamma_i = -\overline{\gamma}\ln(U_i), \quad U_i \sim \mathrm{Uniform}(0,1).
\end{equation}

\textbf{Cosine-annealed thresholds.} To protect gradient quality during early training, we define:
\begin{align}
    \gamma_{\mathrm{high}}(t) &= \gamma_h^0\!\cdot\!\left[1 + \tfrac{1}{2}\cos\!\left(\tfrac{\pi t}{T}\right)\right], \quad \gamma_h^0 = 15\,\text{dB}, \\
    \gamma_{\mathrm{low}}(t) &= \gamma_l^0\!\cdot\!\left[1 + \tfrac{1}{2}\cos\!\left(\tfrac{\pi t}{T}\right)\right], \quad \gamma_l^0 = 5\,\text{dB}.
\end{align}
At $t=0$: $[\gamma_{\mathrm{high}}, \gamma_{\mathrm{low}}] = [22.5, 7.5]$\,dB (conservative). At $t=T$: $[7.5, 2.5]$\,dB (aggressive).

\textbf{Three-level payload strategy.}
\begin{equation}
    \text{payload}_i = \begin{cases}
        \left(\mathbf{B}_i,\, \mathbf{h}_1^{\mathrm{fp32}},\, \mathbf{h}_2^{\mathrm{fp32}}\right) & \gamma_i > \gamma_{\mathrm{high}}(t), \\[4pt]
        \left(\mathbf{B}_i,\, \mathbf{h}_1^{\mathrm{int8}},\, \mathbf{h}_2^{\mathrm{int8}},\, s_1, s_2\right) & \gamma_l < \gamma_i \leq \gamma_h, \\[4pt]
        \mathbf{B}_i & \gamma_i \leq \gamma_{\mathrm{low}}(t).
    \end{cases}
    \label{eq:payload}
\end{equation}
For INT8 quantization of $\mathbf{h}_k$:
\begin{equation}
    s_k = \frac{\max|\mathbf{h}_k|}{127}, \quad \mathbf{h}_k^q = \left\lfloor\frac{\mathbf{h}_k}{s_k}\right\rceil_{\![-127,127]}^{\mathrm{int8}}.
\end{equation}
The quantization error satisfies $\|\mathbf{h}_k - s_k \mathbf{h}_k^q\|_\infty \leq s_k/2$.

\textbf{Expected communication.} At $\overline{\gamma}=10$\,dB and $t=T/2$ ($\gamma_h=15$\,dB, $\gamma_l=5$\,dB):
\begin{align}
    P(\mathrm{FP32}) &= e^{-\gamma_h/\overline{\gamma}} \approx 4.1\%, \\
    P(\mathrm{INT8}) &= e^{-\gamma_l/\overline{\gamma}} - e^{-\gamma_h/\overline{\gamma}} \approx 68.6\%, \\
    P(\mathrm{bin}) &= 1 - e^{-\gamma_l/\overline{\gamma}} \approx 27.3\%.
\end{align}

\subsection{NIA-CVA Aggregation}

Standard FedAvg aggregates with data-proportional weights $w_i = n_i/n$. NIA-CVA introduces error- and channel-aware reweighting:
\begin{equation}
    \tilde{w}_i = w_i\cdot\psi(\gamma_i, \varepsilon_i), \quad
    \tilde{w}_i \leftarrow \frac{\tilde{w}_i}{\sum_j \tilde{w}_j},
\end{equation}
where
\begin{equation}
    \psi(\gamma_i, \varepsilon_i) = \left(1 - \alpha_\varepsilon \frac{\varepsilon_i}{\varepsilon_{\max}}\right)\!\left(1 - \alpha_b\,\mathbf{1}[\text{binary-only}_i]\right),
\end{equation}
$\varepsilon_i = \|\Delta\mathbf{W}_i - \widehat{\Delta\mathbf{W}}_i\|_F$ is the rank-1 reconstruction error (transmitted in payload metadata), $\alpha_\varepsilon = 0.3$, and $\alpha_b = 0.5$. The aggregation variance under this scheme satisfies:
\begin{equation}
    \mathbb{E}\!\left[\left\|\mathbf{G}_{\mathrm{agg}} - \mathbf{G}_{\mathrm{true}}\right\|^2\right] = \sum_i \tilde{w}_i^2\,\mathrm{Var}(\mathbf{G}_i),
\end{equation}
which is minimized by assigning lower $\tilde{w}_i$ to high-$\mathrm{Var}(\mathbf{G}_i)$ clients.

\subsection{Adaptive Frequency Calibration (AFC)}

Binary quantization $\mathbf{B} = \mathrm{sign}(\Delta\mathbf{W})$ amplifies the dominant singular direction, introducing systematic spectral bias:
\begin{equation}
    \mathbb{E}\!\left[\mathbf{B} \odot \left(\mathbf{h}_1\mathbf{h}_2^\top\right)\right] \neq \Delta\mathbf{W}.
\end{equation}
AFC corrects this using inter-round gradient consistency:
\begin{align}
    \beta(t) &= \beta_0 \cdot \max\!\left(0,\; \frac{\langle \mathbf{G}_{\mathrm{agg}},\, \mathbf{G}_{\mathrm{prev}}\rangle_F}{\|\mathbf{G}_{\mathrm{agg}}\|_F \|\mathbf{G}_{\mathrm{prev}}\|_F}\right), \\
    \mathbf{G}_{\mathrm{cal}} &= \mathbf{G}_{\mathrm{agg}} + \beta(t)\left(\mathbf{G}_{\mathrm{agg}} - \mathbf{G}_{\mathrm{prev}}\right),
\end{align}
with $\beta_0 = 0.05$. When consecutive global gradients are well-aligned (consistent signal), AFC amplifies the update to recover suppressed frequency components; when they diverge (noisy round), $\beta \to 0$.

%======================================================
\section{Theoretical Analysis}
%======================================================

\subsection{Convergence Bound}

\textbf{Assumptions.} (A1) Each $F_i$ is $L$-smooth. (A2) Stochastic gradient variance is bounded: $\mathbb{E}[\|\nabla F_i(\mathbf{w},\xi) - \nabla F_i(\mathbf{w})\|^2] \leq \sigma^2$. (A3) Bounded gradient dissimilarity: $\mathbb{E}[\|\nabla F_i(\mathbf{w}) - \nabla F(\mathbf{w})\|^2] \leq \Gamma^2$.

\begin{theorem}[RC-FedBiT Convergence]
Under (A1)--(A3), with learning rate $\eta = \mathcal{O}(1/\sqrt{T})$:
\begin{equation}
    \frac{1}{T}\sum_{t=0}^{T-1}\mathbb{E}\!\left[\left\|\nabla F(\mathbf{W}^{(t)})\right\|^2\right]
    \leq \mathcal{O}\!\left(\frac{\sigma}{\sqrt{KT}}\right) + \delta_{\mathrm{comp}} + \delta_{\mathrm{ch}} + \mathcal{O}\!\left(\eta^2 L E^2 \Gamma^2\right),
\end{equation}
where
\begin{align}
    \delta_{\mathrm{comp}} &= \frac{L\eta\,\mathbb{E}[\varepsilon^2]}{Kmn}, \\
    \delta_{\mathrm{ch}} &= \frac{L\eta\, P_{\mathrm{bin}}(t)\,\sigma_{\mathrm{bin}}^2}{K}, \quad P_{\mathrm{bin}}(t) = 1 - e^{-\gamma_{\mathrm{low}}(t)/\overline{\gamma}}.
\end{align}
\end{theorem}

As $\overline{\gamma}\to\infty$, $P_{\mathrm{bin}}(t)\to 0$ and the channel noise term $\delta_{\mathrm{ch}}$ vanishes. Cosine annealing of $\gamma_{\mathrm{low}}(t)$ further reduces $P_{\mathrm{bin}}(t)$ as training progresses.

%======================================================
\section{Experiments}
%======================================================

\subsection{Setup}

Table~\ref{tab:setup} summarizes the experimental configuration. All experiments use ViT-Tiny~\cite{dosovitskiy2021image} (patch size 4, image size 32, $\sim\!5.7$M parameters) on CIFAR-10~\cite{krizhevsky2009learning}. Client data is partitioned via the Dirichlet distribution with $\alpha=0.5$.

\begin{table}[h]
\centering
\caption{Experimental Configuration}
\label{tab:setup}
\begin{tabular}{lc|lc}
\toprule
\textbf{Param} & \textbf{Value} & \textbf{Param} & \textbf{Value} \\
\midrule
Dataset & CIFAR-10 & Model & ViT-Tiny \\
Clients $N$ & 100 & Participation $\rho$ & 0.1 \\
Data split & Dir($\alpha$=0.5) & Local epochs $E$ & 5 \\
Learning rate $\eta$ & 0.01 & Rounds $T$ & 100 \\
Mean SNR $\overline{\gamma}$ & 10 dB & GPU & RTX 4090 \\
$\gamma_h^0$ / $\gamma_l^0$ & 15 / 5 dB & $\beta_0$ & 0.05 \\
\bottomrule
\end{tabular}
\end{table}

\subsection{Main Results}

Table~\ref{tab:main} reports final test accuracy and communication cost at round $T=100$. RC-FedBiT achieves accuracy within $0.25$ pp of full-precision FedAvg while using only $\sim\!4$--$14\%$ of its communication budget, representing a $7$--$25\times$ reduction.

\begin{table}[h]
\centering
\caption{Main Comparison on CIFAR-10 ($T=100$ rounds)}
\label{tab:main}
\begin{tabular}{lccc}
\toprule
\textbf{Method} & \textbf{Acc (\%)} & \textbf{Rel. Cost} & $\boldsymbol{\Delta}$ \textbf{vs FedAvg} \\
\midrule
FedAvg~\cite{mcmahan2017communication} & 26.54 & 100\% & --- \\
\textbf{RC-FedBiT (ours)} & \textbf{26.29} & $\sim$4--14\% & $-$0.25 pp \\
QSGD~\cite{alistarh2017qsgd} & 28.06 & $\sim$25\% & $+$1.52 pp \\
SignSGD~\cite{bernstein2018signsgd} & 23.26 & $\sim$3\% & $-$3.28 pp \\
PowerSGD~\cite{vogels2019powersgd} & 24.10 & $\sim$30\% & $-$2.44 pp \\
\bottomrule
\end{tabular}
\end{table}

RC-FedBiT outperforms all low-communication baselines. SignSGD suffers a $3.28$ pp drop due to complete magnitude suppression. PowerSGD's $2.44$ pp gap reflects its channel-oblivious design. QSGD achieves the highest accuracy ($28.06\%$) via unbiased stochastic quantization, but at $\sim\!25\%$ cost with no channel robustness guarantee. RC-FedBiT's $1.77$ pp gap versus QSGD represents the accuracy trade-off for channel robustness, which is expected to reverse at low SNR conditions (experiments pending).

\subsection{Ablation Study}

Table~\ref{tab:ablation} ablates each RC-FedBiT component. All variants use $T=100$ rounds on CIFAR-10.

\begin{table}[h]
\centering
\caption{Ablation Study ($T=100$ rounds, CIFAR-10)}
\label{tab:ablation}
\begin{tabular}{llcc}
\toprule
\textbf{Variant} & \textbf{Active Components} & \textbf{Acc (\%)} & \textbf{Status} \\
\midrule
A0 (baseline) & Binary only & ${\sim}22.5^*$ & Est. \\
A1 & $+$ Rank-1 & ${\sim}25.0^*$ & Est. \\
A2 (ca\_rs) & $+$ Rank-1 $+$ CA-RS & 27.13 & Measured \\
A3 (nia\_cva) & $+$ CA-RS $+$ NIA-CVA & 27.13 & Measured \\
A4 (full) & All four components & 27.01 & Measured \\
\midrule
FedAvg & Full precision (reference) & 26.54 & Measured \\
\bottomrule
\end{tabular}
\small{$^*$Estimated from partial runs; full results pending.}
\end{table}

\textbf{CA-RS} (A1$\to$A2, estimated $+2.1$ pp): The dominant component. By matching precision to channel capacity, CA-RS prevents low-SNR clients from corrupting the aggregation with noisy FP32 transmissions. Both A2 and A3 exceed FedAvg ($27.13\%$ vs $26.54\%$), demonstrating that channel-adaptive compression can \emph{outperform} full-precision aggregation in wireless FL.

\textbf{NIA-CVA} (A2$\to$A3, $0.00$ pp): No measurable gain in the IID Dirichlet($\alpha{=}0.5$) setting. The benefit of NIA-CVA is expected to emerge under high data heterogeneity ($\alpha < 0.1$), where gradient variance across clients is large. Non-IID experiments are ongoing.

\textbf{AFC} (A3$\to$A4, $-0.12$ pp): A small regression over $100$ rounds. AFC may require longer training ($T > 200$) for spectral correction benefits to manifest, or the current $\beta_0 = 0.05$ schedule may be too aggressive.

%======================================================
\section{Related Work}
%======================================================

\textbf{Gradient compression.}
SignSGD~\cite{bernstein2018signsgd} transmits only gradient signs, reducing cost dramatically but losing magnitude information. QSGD~\cite{alistarh2017qsgd} uses stochastic uniform quantization with unbiasedness guarantees. PowerSGD~\cite{vogels2019powersgd} approximates gradients with rank-$r$ factors for flexible compression. 1-bit Adam~\cite{tang20211} combines momentum factors with error feedback for further compression. RC-FedBiT differs from all of these by explicitly incorporating channel quality into compression decisions.

\textbf{Communication-efficient FL.}
FedProx~\cite{li2020federated} adds a proximal term for heterogeneous settings. SCAFFOLD~\cite{karimireddy2020scaffold} uses control variates to reduce client drift. None of these works address wireless channel-adaptive bit-width allocation.

\textbf{Channel-aware FL.}
Recent works~\cite{amiri2020machine,zhu2019broadband} study over-the-air computation for FL in noisy channels, but typically assume analog transmission. RC-FedBiT instead operates on digital bit-width adaptation, compatible with standard wireless protocols.

%======================================================
\section{Conclusion}
%======================================================

We presented RC-FedBiT, a rate-controlled federated learning framework that adapts gradient bit-width to instantaneous wireless channel conditions. By combining Rank-1 compression, CA-RS, NIA-CVA, and AFC, RC-FedBiT achieves accuracy within $0.25$ pp of full-precision FedAvg at $\sim\!25\%$ communication cost. Ablation variants surpass FedAvg ($27.13\%$ vs $26.54\%$), demonstrating that channel-adaptive compression can outperform full-precision aggregation in wireless settings.

Future directions include: (1)~completing channel SNR and non-IID experiments; (2)~formal convergence analysis with explicit channel-dependent error bounds; (3)~extension to personalized FL with client-specific spectral calibration; (4)~hardware-efficient INT4 and binary codecs for embedded FL devices.

%======================================================
\bibliographystyle{IEEEtran}
\begin{thebibliography}{10}

\bibitem{mcmahan2017communication}
B.~McMahan, E.~Moore, D.~Ramage, S.~Hampson, and B.~A. y~Arcas,
``Communication-efficient learning of deep networks from decentralized data,''
in \emph{AISTATS}, 2017.

\bibitem{bernstein2018signsgd}
J.~Bernstein, Y.-X. Wang, K.~Azizzadenesheli, and A.~Anandkumar,
``signSGD: Compressed optimisation for non-convex problems,''
in \emph{ICML}, 2018.

\bibitem{alistarh2017qsgd}
D.~Alistarh, D.~Grubic, J.~Li, R.~Tomioka, and M.~Vojnovic,
``QSGD: Communication-efficient SGD via gradient quantization and encoding,''
in \emph{NeurIPS}, 2017.

\bibitem{vogels2019powersgd}
T.~Vogels, S.~P. Karimireddy, and M.~Jaggi,
``PowerSGD: Practical low-rank gradient compression for distributed optimization,''
in \emph{NeurIPS}, 2019.

\bibitem{tang20211}
S.~Tang \emph{et al.},
``1-bit Adam: Communication efficient large-scale training with Adam's convergence speed,''
in \emph{ICML}, 2021.

\bibitem{li2020federated}
T.~Li, A.~K. Sahu, M.~Zaheer, M.~Sanjabi, A.~Smola, and V.~Smith,
``Federated optimization in heterogeneous networks,''
in \emph{MLSys}, 2020.

\bibitem{karimireddy2020scaffold}
S.~P. Karimireddy, S.~Kale, M.~Mohri, S.~Reddi, S.~Stich, and A.~T. Suresh,
``SCAFFOLD: Stochastic controlled averaging for federated learning,''
in \emph{ICML}, 2020.

\bibitem{amiri2020machine}
M.~M. Amiri and D.~G{\"u}nd{\"u}z,
``Machine learning at the wireless edge: Distributed stochastic gradient descent over-the-air,''
\emph{IEEE Trans. Signal Process.}, 2020.

\bibitem{zhu2019broadband}
G.~Zhu, Y.~Wang, and K.~Huang,
``Broadband analog aggregation for low-latency federated edge learning,''
\emph{IEEE Trans. Wireless Commun.}, 2019.

\bibitem{dosovitskiy2021image}
A.~Dosovitskiy \emph{et al.},
``An image is worth 16x16 words: Transformers for image recognition at scale,''
in \emph{ICLR}, 2021.

\bibitem{krizhevsky2009learning}
A.~Krizhevsky,
``Learning multiple layers of features from tiny images,''
Tech. Rep., 2009.

\end{thebibliography}

\end{document}
